% !TeX program = xelatex
\documentclass[tikz,10.5pt]{standalone}
\usepackage[UTF8,scheme=plain,fontset=windows]{ctex}
\usepackage{tikz}
\usetikzlibrary{arrows.meta,calc,positioning,fit}
\definecolor{nodeBlue}{RGB}{58,139,198}
\definecolor{dataGreen}{RGB}{68,157,92}
\definecolor{commRed}{RGB}{198,74,78}
\definecolor{ctrlGray}{RGB}{111,123,138}
\definecolor{softBG}{RGB}{247,249,251}
\definecolor{ink}{RGB}{36,42,50}
\begin{document}
\begin{tikzpicture}[
  >=Latex,
  font=\small,
  panel/.style={draw=ink!22, fill=softBG, rounded corners=5pt, minimum height=5.15cm},
  qpu/.style={circle, draw=nodeBlue!75!black, fill=nodeBlue!18, minimum size=7.5mm},
  resource/.style={draw=nodeBlue!65!black, line width=0.65pt},
  smallbox/.style={draw=ink!35, fill=white, rounded corners=3pt, align=center, font=\scriptsize, inner sep=3pt},
  flowbox/.style={draw=nodeBlue!65!black, fill=nodeBlue!8, rounded corners=3pt, align=center, font=\scriptsize, minimum height=7mm, inner sep=3pt},
  arrow/.style={-{Latex[length=2mm]}, line width=0.65pt, draw=ink!60}
]

% Panel (a)
\node[panel, minimum width=5.55cm] (pa) at (0,0) {};
\node[anchor=north west, font=\small, text=ink] at ($(pa.north west)+(0.15,-0.12)$) {(a) 模块网络};
\node[qpu] (a1) at (-1.85,1.05) {QPU};
\node[qpu] (a2) at (0.05,1.45) {};
\node[qpu] (a3) at (1.85,0.75) {};
\node[qpu] (a4) at (-1.3,-1.05) {};
\node[qpu] (a5) at (1.05,-1.2) {};
\draw[resource] (a1) -- (a2);
\draw[resource] (a2) -- (a3);
\draw[resource] (a1) -- (a4);
\draw[resource] (a4) -- (a5);
\draw[resource] (a5) -- (a3);
\draw[resource] (a2) -- (a5);
\node[smallbox, fill=nodeBlue!6] at (0,-2.03) {可宣告 Bell 对\\$F_{\mathrm{Bell}},\,R_{\mathrm{link}},\,\tau_h,\,m$};
\node[font=\scriptsize, text=ink!70, align=center] at (0,0.05) {网络边 =\\可调度纠缠资源};

% Panel (b)
\node[panel, minimum width=6.05cm] (pb) at (6.1,0) {};
\node[anchor=north west, font=\small, text=ink] at ($(pb.north west)+(0.15,-0.12)$) {(b) QPU 节点内部};
\node[draw=ink!30, fill=white, rounded corners=5pt, minimum width=4.9cm, minimum height=3.55cm] (qpubox) at (6.1,-0.05) {};
\foreach \x in {-0.9,-0.45,0,0.45,0.9} {
  \foreach \y in {0.35,0.7,1.05} {
    \fill[dataGreen] ($(6.1,0)+(\x,\y)$) circle (1.5pt);
  }
}
\foreach \x in {-1.25,1.25} {
  \fill[commRed] ($(6.1,0)+(\x,-0.35)$) circle (2.2pt);
  \draw[commRed!70, line width=0.55pt] ($(6.1,0)+(\x,-0.35)$) -- ($(6.1,0)+(\x,-0.9)$);
}
\node[smallbox, fill=dataGreen!8, anchor=west] at (6.55,1.08) {数据比特\\本地门 / 存储 / 读出\\综合征循环};
\node[smallbox, fill=commRed!8, anchor=west] at (6.55,-0.55) {通信比特\\原子--光子接口\\远程纠缠收发};
\node[smallbox, fill=ctrlGray!10, minimum width=2.8cm] at (5.35,-1.65) {Pauli frame\\路由选择\\成功记录};
\draw[arrow] (5.35,-1.2) -- (5.35,-0.65);
\draw[arrow] (6.75,-1.2) -- (7.12,-0.82);

% Panel (c)
\node[panel, minimum width=6.75cm] (pc) at (12.65,0) {};
\node[anchor=north west, font=\small, text=ink] at ($(pc.north west)+(0.15,-0.12)$) {(c) 远程资源调用链};
\node[flowbox] (r1) at (10.05,0.75) {可宣告\\Bell 对};
\node[flowbox] (r2) at (11.35,0.75) {纯化};
\node[flowbox] (r3) at (12.7,0.75) {纠缠\\交换};
\node[flowbox] (r4) at (14.05,0.75) {门传送};
\node[flowbox, fill=dataGreen!9, draw=dataGreen!70!black] (r5) at (15.35,0.75) {跨模块门\\逻辑操作};
\draw[arrow] (r1) -- (r2);
\draw[arrow] (r2) -- (r3);
\draw[arrow] (r3) -- (r4);
\draw[arrow] (r4) -- (r5);
\draw[commRed!65, line width=0.65pt] (10.05,-0.55) -- (15.35,-0.55);
\foreach \x in {10.05,12.7,15.35} {
  \fill[commRed] (\x,-0.55) circle (2pt);
}
\node[smallbox, fill=commRed!8] at (11.35,-1.22) {heralded success\\成功记录 $m$};
\node[font=\scriptsize, align=center, text=ink!72] at (14.0,-1.22) {资源进入\\上层调度};
\draw[arrow, draw=commRed!70!black] (11.35,-0.92) -- (11.35,-0.58);
\draw[arrow, draw=dataGreen!70!black] (14.0,-0.95) -- (14.0,-0.58);

\end{tikzpicture}
\end{document}
